\section{Background}\label{sec:background}

\subsection{Test oracles and derived oracles}
A test oracle decides whether an observed behaviour is correct \cite{barr2015oracle}. Barr et al.\ distinguish specified oracles (from a formal specification), derived oracles (from other artefacts such as documentation, earlier versions or properties of the program), implicit oracles (crashes, exceptions) and human oracles. Metamorphic testing \cite{chen2018metamorphic,segura2016survey} is the best-known family of derived oracles: instead of the correct output for one input, it checks a relation between outputs of related executions. Metamorphic and differential oracles are well established for scientific software \cite{kanewala2014scientific}, autonomous-driving software \cite{zhou2019driverless} and graphics drivers \cite{donaldson2017shader}, where exact expected outputs are unavailable. Recent work generates oracles or assertions with neural models \cite{dinella2022toga,hossain2023neural} and discusses the risks of letting an LLM act as the oracle for code that an LLM wrote \cite{molina2025oracle}. For LLM-generated code in particular, weak test suites accept many incorrect programs \cite{liu2023evalplus}.

\subsection{Kinematics of planar mechanisms}
All mechanisms in this study are planar and driven by a single input, a crank turning or swinging at constant angular speed. The position of every joint follows from the input angle by solving loop-closure equations. For the linkages used here, each unknown joint is the intersection of two circles centred on joints already known. Each intersection has two solutions, the \emph{assembly branches}, and a correct simulator must stay on one branch throughout the cycle. Whether a solution exists for every input angle (the linkage \emph{assembles}) depends on the link lengths. The Grashof condition, for example, decides whether a four-bar's crank can turn fully. Beyond assembly, many mechanisms work only because of specific \emph{design knowledge}:
\begin{itemize}
\item \emph{published dimensions}, such as Jansen's eleven link lengths \cite{jansen_strandbeest} or the proportions of Chebyshev's lambda linkage, crank $a$, ground $2a$, rocker $2.5a$, coupler $2.5a$ extended by a further $2.5a$ to the tracing point \cite{wiki_lambda};
\item \emph{relations}, such as the Peaucellier--Lipkin conditions (equal rhombus sides, equal long arms, crank length equal to the distance between the fixed pivots), under which the output point moves on an exact straight line;
\item \emph{assembly conditions}, such as the requirement that the sum of sun and ring teeth of a planetary set be divisible by the number of equally spaced planets.
\end{itemize}
Errors in design knowledge are particularly treacherous. The mechanism still assembles and moves smoothly, but it does not do what it is for.

\subsection{Generated simulators and their acceptance}
We consider simulators written as one self-contained HTML page using \texttt{three.js}, as produced by the \emph{teardown} skill for Claude Code \cite{teardown}. The skill's own review loop has the model look at sheets of rendered frames and correct what it sees. The same pattern appears in research systems that use a VLM as critic of generated 3D or physics code \cite{le2025articulate,alrashedy2025cad,shende2026perceptual,gao2026lam}. In the terms of simulation V\&V \cite{sargent2013vv,balci1997vva}, this is the \emph{animation} technique: watching the model run to judge its behaviour. It is useful for gross errors, but it cannot reveal quantities that are not visible and is prone to rationalising small discrepancies. The comparison of a model's output with an analytic solution, the other classical technique, has so far not been applied to generated simulators. Making it possible requires access to the simulator's internal state at arbitrary times, which motivates the minimal output contract of Section~\ref{sec:approach}.
